\subsection{Special case of a covering consisting of two parts}

Let X be a path-connected topological space, and let B and C be nonempty path-connected subsets of X. Moreover, assume one of the following two hypotheses:

\begin{enumerate}
\def\labelenumi{(\roman{enumi})}
\item
  The interiors of the sets B and C cover X;
\item
  The sets B and C are closed in X, their union is equal to X, the spaces X, B, C are delaceable, and the space B ∩ C is locally path-connected.
\end{enumerate}

Under these hypotheses, the canonical map from the sum space of the family (B, C) to the space X satisfies property (VK) (cf. IV, p. 409).

Set A = B ∩ C. Since the space X is connected, the set A is nonempty. Let a be a point of A; let \(j_{0}\) be the path-connected component of a in A. For every path-connected component j of A distinct from \(j_{0}\) , choose a point \(a_{j}\) of j, the class \(\beta_{j}\) of a path in B and the class \(\gamma_{j}\) of a path in C, both joining \(a_{j}\) to a; denote \(\varphi_{j}\colon\pi_{1}(A,a_{j})\to\pi_{1}(B,a)\) and \(\psi_{j}\colon\pi_{1}(A,a_{j})\to\pi_{1}(C,a)\) the group homomorphisms defined by

\[
\varphi_ {j} (v) = \beta_ {j} ^ {- 1} v \beta_ {j} \quad \text {and} \quad \psi_ {j} (v) = \gamma_ {j} ^ {- 1} v \gamma_ {j}
\]

for \(v \in \pi_1(A, a_j)\) . We also denote by \(\varphi_0\) and \(\psi_0\) the homomorphisms from \(\pi_1(A, a)\) into \(\pi_1(B, a)\) and \(\pi_1(C, a)\) respectively deduced from the canonical injections. Denote by \(\iota_B\) and \(\iota_C\) the homomorphisms from \(\pi_1(B, a)\) and \(\pi_1(C, a)\) respectively into \(\pi_1(X, a)\) deduced from the canonical injections. Finally, let \(\mu\) be the unique homomorphism from the free group \(F(\pi_0(A) - \{j_0\})\) into \(\pi_1(X, a)\) such that \(\mu(j) = \beta_j^{-1}\gamma_j\) for all \(j \in \pi_0(A) - \{j_0\}\) .

PROPOSITION 2. --- There exists a unique group homomorphism

\[
\mathsf {M} \colon \pi_ {1} (\mathrm{B}, a) * \pi_ {1} (\mathrm{C}, a) * \mathrm{F} (\pi_ {0} (\mathrm{A}) - \{j _ {0} \}) \to \pi_ {1} (\mathrm{X}, a)
\]

which coincides with \(\iota_{\mathrm{B}}\) on \(\pi_1(\mathrm{B},a)\) , \(\iota_{\mathrm{C}}\) on \(\pi_1(\mathrm{C},a)\) and \(\mu\) on \(\mathrm{F}(\pi_0(\mathrm{A}) - \{j_0\})\) . This homomorphism is surjective and its kernel is the smallest distinguished subgroup containing the elements

\[
\varphi_ {j} (v) j \psi_ {j} (v) ^ {- 1} j ^ {- 1},
\]

for \(j\in \pi_0(\mathrm{A}) - \{j_0\}\) and \(v\in \pi_1(\mathrm{A},a_j)\) , and the elements

\[
\varphi_ {0} (v) \psi_ {0} (v) ^ {- 1}
\]

for \(v \in \pi_1(\mathrm{A}, a)\) .

The frame \(\Gamma\) of the covering of X defined by the family (B,C) has two vertices \(b\) and \(c\) corresponding to the two sets B and C. The set of its arrows is \(\pi_0(\mathrm{A})\) ; they join the vertices \(b\) and \(c\) . The graph associated with the subquiver of \(\Gamma\) whose unique arrow is \(j_0\) is a maximal tree of \(\widetilde{\Gamma}\) . The proposition then follows from IV, p. 412, prop. 1.

Example. --- For \(n \geqslant 1\) the sphere \(S_{n}\) is the union of two closed hemispheres, homeomorphic to the closed ball \(B_{n-1}\) (TG, VI, p. 12), and whose intersection is identified with the sphere \(S_{n-1}\) . For \(n \geqslant 2\) the sphere \(S_{n-1}\) is path-connected; one deduces that the Poincaré group of \(S_{n}\) is trivial (cf. I, p. 127, example 3).

The sphere \(S_{0}\) has two path-connected components; one thus recovers that the Poincaré group of the circle \(S_{1}\) is isomorphic to a free group with one generator. More precisely, let B and C be the intersections of \(S_{1}\) with the half-planes of equations \(y \geqslant 0\) and \(y \leqslant 0\) in the plane \(R^{2}\) . Set \(a = (1,0)\) , \(a' = (-1,0)\) ; we have \(B \cap C = \{a, a'\}\) ; its connected components are \(j_{0} = \{a\}\) and \(j = \{a'\}\) . Let \(\beta\) be the class of the path \(t \mapsto e^{\pi it}\) in C; if one identifies C with \(R^{2}\) , it joins a to \(a'\) in B. Likewise, let \(\gamma\) be the class of the path \(t \mapsto e^{-\pi it}\) connecting a to \(a'\) in C. The path \(\beta\gamma^{-1}\) is a loop at a, given by \(t \mapsto e^{2\pi it}\) . By proposition 2, its class generates the group \(\pi_{1}(\mathbf{S}_{1}, a)\) .

COROLLARY 1. --- The homomorphism \(\mu\) is injective. More precisely, there exists a retraction associated with \(\mu\) which is a group homomorphism.

Let \(\rho\) be the unique homomorphism from \(\pi_1(\mathrm{B},a)*\pi_1(\mathrm{C},a)*\mathrm{F}(\pi_0(\mathrm{A})-\{j_0\})\) into \(\mathrm{F}(\pi_0(\mathrm{A}) - \{j_0\})\) which induces the trivial homomorphism on \(\pi_1(\mathrm{B},a)\) and \(\pi_1(\mathrm{C},a)\) and the identity on \(\mathrm{F}(\pi_0(\mathrm{A}) - \{j_0\})\) . Let N be the kernel of the homomorphism M. By proposition 2, \(\rho (\mathbf{N})\) is reduced to the identity element. There therefore exists a unique homomorphism r from \(\pi_1(\mathrm{X},a)\) into \(\mathrm{F}(\pi_0(\mathrm{A}) - \{j_0\})\) such that \(\rho = r\circ M\) . For all \(v\in \mathrm{F}(\pi_0(\mathrm{A}) - \{j_0\})\) , we have \(\mathsf{M}(v) = \mu (v)\) , therefore \(r\circ \mu\) is the identity homomorphism. The corollary follows.

COROLLARY 2. --- If the group \(\pi_{1}(X,a)\) is trivial, the set \(A = B \cap C\) is path-connected; if it is commutative, the set A has at most two path-connected components.

Indeed, if S is a set, the free group F(S) is trivial only if S is empty and is commutative only if Card S ≤ 1.

COROLLARY 3. --- If the groups \(\pi_1(\mathrm{X},a)\) and \(\pi_1(\mathrm{A},a)\) are trivial, so are the groups \(\pi_1(\mathrm{B},a)\) and \(\pi_1(\mathrm{C},a)\) .

By corollary 2, the set A is path-connected. The group \(\pi_1(\mathrm{X},a)\) is therefore isomorphic to the free product of the groups \(\pi_{1}(B,a)\) and \(\pi_{1}(C,a)\) . It contains in particular subgroups isomorphic to the groups \(\pi_{1}(B,a)\) and \(\pi_{1}(C,a)\) (A, I, p. 83). These two groups are therefore trivial if \(\pi_{1}(X,a)\) is trivial.

